\section{Worked case: transcriptional memory through chromatin succession}
\label{sec:ino1}

\subsection{The system}

Budding yeast cells previously starved of inositol reactivate the \textit{INO1}
gene faster on a second starvation than cells without that history. The memory
requires a cis-acting DNA element (the memory recruitment sequence, MRS), the
transcription factor Sfl1 bound to it, interaction with the nuclear pore
complex, incorporation of H2A.Z, and a memory-specific dimethylation of histone
H3 lysine 4 (H3K4me2) \citep{LightEtAl2010INO1,SumpEtAl2022INO1Memory}. The mark
is written by a form of COMPASS lacking the Spp1 subunit and requires SET3C, a
complex that binds H3K4me2 and interacts physically with the writer; the
authors propose this writer--reader coupling as the mechanism for
re-incorporating the mark after DNA replication \citep{SumpEtAl2022INO1Memory}.

Sump and colleagues separated establishment from inheritance by inducible
degradation of Sfl1. When Sfl1 was degraded after memory had been established,
peripheral localization was lost but H3K4me2 remained at the promoter for
roughly three to four generations; degrading Set3 together with Sfl1 caused
rapid loss of the mark \citep{SumpEtAl2022INO1Memory}. The authors conclude
that memory duration is governed by the number of divisions or replication
cycles, reflecting dilution, degradation, or imperfect fidelity of inheritance.

\subsection{Mapping onto the criterion}

Each replication fork distributes parental nucleosomes between the daughter
chromatids and incorporates new ones, so the physical realizations
$c_0,c_1,c_2,\ldots$ of the mark are progressively disjoint in material
composition. This is the situation the succession clause of
Section~\ref{sec:criterion} was written for. Table~\ref{tab:ino1} gives the
mapping.

\begin{table}[ht]
\centering
\small
\caption{The persistent-carrier criterion applied to \textit{INO1}
transcriptional memory. ``Shown'' refers to published results
\citep{LightEtAl2010INO1,SumpEtAl2022INO1Memory}; ``untested'' marks components
of the criterion that no published experiment yet addresses.}
\label{tab:ino1}
\begin{tabularx}{\linewidth}{>{\raggedright\arraybackslash}p{0.2\linewidth}
  >{\raggedright\arraybackslash}X
  >{\raggedright\arraybackslash}p{0.13\linewidth}}
\toprule
Element & Instantiation & Status \\
\midrule
History $m$ vs.\ $m_0$ & Prior inositol starvation vs.\ never-starved cells in the
same current medium & Shown \\
Future family $\Efamily$ & Distribution of \textit{INO1} re-induction times
after re-starvation, indexed by the delay $\pi\in\Pi$ & Shown (faster
reactivation) \\
Constraint type & Soft: naive cells also induce, so the support is unchanged
and the weights shift; Eq.~\eqref{eq:icp-constraint-divergence} applies
directly & Inferred \\
Formation & Sfl1 and pore interaction required to establish the mark & Shown \\
Carrier intervention & Removing writer/reader components (COMPASS, SET3C)
abolishes the mark and faster reactivation; matched shams required & Shown,
shams partial \\
Succession & H3K4me2 re-incorporated onto new nucleosomes for 3--4 generations
after initiator removal & Shown for the mark \\
Role equivalence $\sim$ & Successor cells whose re-induction distribution lies
within a declared tolerance of the memory distribution & Untested across
succession \\
Role exhaustion & Mark lost after finitely many divisions while the lineage
lives: $\Rset_C=\varnothing$ for the role, not the substrate & Shown for the
mark \\
\bottomrule
\end{tabularx}
\end{table}

\subsection{What the case shows and does not}

The case passes formation, sufficiency, and carrier intervention for the mark,
and it demonstrates succession of the \emph{record} through several rounds of
replication-coupled carrier replacement. It does not yet demonstrate
succession of the \emph{role}. Sfl1 is itself required to express the faster
reactivation, so while Sfl1 is absent the future-response family cannot be
assayed: the inherited mark is a latent record whose functional reading has not
been measured after succession. This distinction between persistence of a
record and persistence of a role is exactly the one ICP exists to enforce, and
the published literature does not need ICP to draw it. It does, however, make
the missing experiment precise.

\paragraph{Prospective test.} Establish memory, degrade Sfl1, allow
$k=0,1,\ldots,5$ generations of latent inheritance, restore Sfl1, and measure
re-induction after re-starvation. ICP's succession clause predicts that the
divergence $\mathfrak C_k$ between memory and naive re-induction-time
distributions decays with $k$ at the per-generation rate of H3K4me2 retention at
the locus, and that degrading Set3 only during the latent interval abolishes
$\mathfrak C_k$ while a matched sham degradation does not. If restored Sfl1
re-establishes fast reactivation independently of the inherited mark, the
role was not carried by the successor record and the case fails the
succession gate.

\paragraph{The ordinary rival.} Nucleosome-modification models with recruited
positive feedback already predict that the stability of a chromatin memory is
set by the balance between feedback-driven re-modification and
replication-coupled dilution, with bistable, effectively heritable states above
a feedback threshold \citep{DoddEtAl2007EpigeneticMemory}. A complete chromatin
model of this kind, with the writer and reader machinery included, fixes the
future family in the sense of Eq.~\eqref{eq:complete-state-null}. We expect
it to predict the outcome of the prospective test, and we do not claim that ICP
explains \textit{INO1} memory better than chromatin dynamics does.

The case is therefore a positive control for the criterion rather than
evidence for a new mechanism. A criterion that failed to classify a known,
mechanistically understood memory correctly would be disqualified; one that
classifies it correctly, separates record from role, and specifies the
untested experiment has earned the right to be applied where the mechanism is
not known. The same table can be filled in for CRISPR spacers, prion
conformational templates, and artificial systems, which makes them comparable
as instances of one question.
